\documentclass{article}
\usepackage{amsmath,amssymb,amsthm}
\usepackage{algorithm,algorithmic}
\usepackage{hyperref}
\newtheorem{theorem}{Theorem}
\newtheorem{lemma}{Lemma}
\newtheorem{assumption}{Assumption}
\newcommand{\EE}{\mathbb{E}}
\newcommand{\norm}[1]{\left\|#1\right\|}
\newcommand{\R}{\mathbb{R}}

\begin{document}

\section*{Clipped Momentum Gradient Descent (CMGD)}

\section*{Mathematical Formulation}
Let $f:\R^{d}\rightarrow\R$ be differentiable.  
At iteration $t\ge 0$ we compute the (possibly stochastic) gradient
\[
g_{t}= \nabla f(x_{t})\quad\text{(or an unbiased estimator).}
\]

\paragraph{Gradient clipping.}
Define the clipping operator with threshold $C>0$:
\[
\tilde g_{t}= \frac{g_{t}}{\displaystyle \max\Bigl(1,\frac{\norm{g_{t}}}{C}\Bigr)}
=
\begin{cases}
g_{t}, & \norm{g_{t}}\le C,\\[4pt]
C\,\dfrac{g_{t}}{\norm{g_{t}}}, & \norm{g_{t}}> C .
\end{cases}
\tag{1}
\]

\paragraph{Momentum on the clipped gradients.}
Let $\beta\in[0,1)$ be the momentum coefficient.  Initialise $m_{-1}=0$ and update
\[
m_{t}= \beta\,m_{t-1}+(1-\beta)\,\tilde g_{t}. \tag{2}
\]

\paragraph{Parameter update.}
With learning rate $\eta>0$,
\[
x_{t+1}= x_{t}-\eta\, m_{t}. \tag{3}
\]

The algorithm is completely deterministic once the (possibly stochastic) gradients $g_{t}$ are fixed.

\section*{Pseudocode}
\begin{algorithm}[H]
\caption{Clipped Momentum Gradient Descent (CMGD)}
\begin{algorithmic}[1]
\REQUIRE initial point $x_{0}\in\R^{d}$, stepsize $\eta>0$, clipping threshold $C>0$, momentum $\beta\in[0,1)$, total iterations $T$.
\STATE $m_{-1}\gets 0$.
\FOR{$t=0,\dots,T-1$}
    \STATE Compute (stochastic) gradient $g_{t}= \nabla f(x_{t})$.
    \STATE \textbf{Clip:}
    \[
    \tilde g_{t}\gets 
    \frac{g_{t}}{\max\bigl(1,\norm{g_{t}}/C\bigr)}.
    \]
    \STATE \textbf{Momentum:}
    \[
    m_{t}\gets \beta\,m_{t-1}+(1-\beta)\,\tilde g_{t}.
    \]
    \STATE \textbf{Update:}
    \[
    x_{t+1}\gets x_{t}-\eta\,m_{t}.
    \]
\ENDFOR
\RETURN $x_{T}$ (or the averaged iterate $\bar x_{T}= \frac1T\sum_{t=0}^{T-1}x_{t}$).
\end{algorithmic}
\end{algorithm}

\section*{Dry Run Example}
Consider the simple quadratic $f(w)=(w-3)^{2}$, $w\in\R$.
\[
\nabla f(w)=2(w-3).
\]
Take $w_{0}=0$, stepsize $\eta=0.1$, clipping threshold $C=1$, momentum $\beta=0.5$.

\begin{center}
\begin{tabular}{c|c|c|c|c}
$t$ & $g_{t}=2(w_{t}-3)$ & $\tilde g_{t}$ (clipped) & $m_{t}$ & $w_{t+1}=w_{t}-\eta m_{t}$ \\ \hline
0 & $-6$ & $-1$ (since $|{-6}|>C$) & $0.5\cdot0+0.5\cdot(-1)=-0.5$ & $0-0.1(-0.5)=0.05$ \\[4pt]
1 & $2(0.05-3)=-5.9$ & $-1$ & $0.5(-0.5)+0.5(-1)=-0.75$ & $0.05-0.1(-0.75)=0.125$ \\[4pt]
2 & $2(0.125-3)=-5.75$ & $-1$ & $0.5(-0.75)+0.5(-1)=-0.875$ & $0.125-0.1(-0.875)=0.2125$ 
\end{tabular}
\end{center}
The objective values after each iteration are
\[
f(w_{0})=9,\; f(w_{1})\approx 8.70,\; f(w_{2})\approx 8.45,\; f(w_{3})\approx 8.20.
\]
Even with aggressive clipping, the method makes progress toward the minimiser $w^{\star}=3$.

\newpage

\section*{Assumptions}
\begin{assumption}[Smoothness]\label{ass:smooth}
$f$ is $L$‑smooth, i.e.
\[
f(y)\le f(x)+\langle\nabla f(x),y-x\rangle+\frac{L}{2}\norm{y-x}^{2},
\qquad\forall x,y\in\R^{d}.
\tag{A1}
\]
\end{assumption}

\begin{assumption}[Bounded variance of stochastic gradients]\label{ass:var}
For every $x$, the stochastic gradient $g(x)$ satisfies
\[
\EE\big[\,g(x)\mid x\,\big]=\nabla f(x),\qquad
\EE\big[\norm{g(x)-\nabla f(x)}^{2}\mid x\big]\le\sigma^{2}.
\tag{A2}
\]
\end{assumption}

\begin{assumption}[Clipping threshold]\label{ass:clip}
The clipping constant $C>0$ is fixed and finite. Consequently
\[
\norm{\tilde g_{t}}\le C\qquad\text{for all }t.
\tag{A3}
\]
\end{assumption}

\begin{assumption}[Momentum parameter]\label{ass:beta}
The momentum coefficient satisfies $0\le\beta<1$.
\tag{A4}
\end{assumption}

All expectations $\EE[\cdot]$ are taken w.r.t. the randomness of the stochastic gradients conditioned on the past iterates.

\section*{Main Convergence Theorem}
\begin{theorem}\label{thm:main}
Under Assumptions \ref{ass:smooth}–\ref{ass:beta}, let the stepsize satisfy
\[
0<\eta\le\frac{1-\beta}{L}.
\tag{T1}
\]
Define the averaged iterate $\bar x_{T}= \frac{1}{T}\sum_{t=0}^{T-1}x_{t}$.  
Then for any $T\ge 1$,
\[
\frac{1}{T}\sum_{t=0}^{T-1}\EE\!\left[\norm{\nabla f(x_{t})}^{2}\right]
\;\le\;
\underbrace{\frac{2\bigl(f(x_{0})-f^{\star}\bigr)}{\eta(1-\beta)T}}_{\text{optimization term}}
\;+\;
\underbrace{2L C^{2}\eta}_{\text{clipping/variance term}}.
\tag{T2}
\]
Consequently, choosing $\eta = \Theta\!\bigl(\tfrac{1}{\sqrt{T}}\bigr)$ yields
\[
\frac{1}{T}\sum_{t=0}^{T-1}\EE\!\left[\norm{\nabla f(x_{t})}^{2}\right]=O\!\bigl(T^{-1/2}\bigr).
\]
\end{theorem}

\section*{Theoretical Properties}
\begin{itemize}
\item \textbf{Stability.} Because clipping guarantees $\norm{\tilde g_{t}}\le C$, the momentum term $m_{t}$ is also uniformly bounded:
\[
\norm{m_{t}}\le C\quad\forall t.
\]
Thus the update (3) cannot produce arbitrarily large jumps, which is crucial for non‑convex landscapes with occasional gradient explosions.

\item \textbf{Hyper‑parameter sensitivity.}
\begin{itemize}
\item The stepsize bound (T1) shows that larger momentum $\beta$ forces a smaller admissible $\eta$.
\item The bound (T2) separates a \emph{deterministic} term (decreases as $1/(\eta T)$) from a \emph{bias} term proportional to $L C^{2}\eta$. The optimal $\eta$ balances these two contributions, giving $\eta\asymp 1/\sqrt{T}$.
\end{itemize}

\item \textbf{Comparison to baselines.}
\begin{itemize}
\item For vanilla SGD (no clipping, $\beta=0$) the standard bound is $O(1/\sqrt{T})$ under the same smoothness and variance assumptions. CMGD attains the same order while providing robustness to gradient explosions.
\item Compared with Adam‑type adaptive methods, CMGD uses only first‑order information and a simple clipping operator, yet its convergence guarantee matches the $O(\log T/T)$ bound of Adam in the non‑convex setting when $C$ is chosen large enough to make clipping inactive.
\end{itemize}
\end{itemize}

\section*{Preliminaries (Definitions and Lemmas)}
\begin{lemma}[Bound on momentum norm]\label{lem:mt}
For all $t\ge0$,
\[
\norm{m_{t}}\le C.
\tag{L1}
\]
\end{lemma}
\begin{proof}
By definition $m_{t}= \beta m_{t-1}+(1-\beta)\tilde g_{t}$ and $\norm{\tilde g_{t}}\le C$ (Assumption \ref{ass:clip}). Using the triangle inequality,
\[
\norm{m_{t}}\le \beta\norm{m_{t-1}}+(1-\beta)C.
\]
Induction on $t$ with $m_{-1}=0$ yields $\norm{m_{t}}\le C$ for every $t$. ∎
\end{proof}

\begin{lemma}[One‑step descent]\label{lem:descent}
For any $t\ge0$,
\[
f(x_{t+1})\le f(x_{t})-\eta(1-\beta)\langle\nabla f(x_{t}),\tilde g_{t}\rangle
+\frac{L\eta^{2}}{2}\norm{m_{t}}^{2}.
\tag{L2}
\]
\end{lemma}
\begin{proof}
From $L$‑smoothness (Assumption \ref{ass:smooth}) and the update $x_{t+1}=x_{t}-\eta m_{t}$,
\[
f(x_{t+1})\le f(x_{t})-\eta\langle\nabla f(x_{t}),m_{t}\rangle+\frac{L\eta^{2}}{2}\norm{m_{t}}^{2}.
\tag{4}
\]
Insert $m_{t}= \beta m_{t-1}+(1-\beta)\tilde g_{t}$ and rearrange:
\[
\langle\nabla f(x_{t}),m_{t}\rangle
= \beta\langle\nabla f(x_{t}),m_{t-1}\rangle+(1-\beta)\langle\nabla f(x_{t}),\tilde g_{t}\rangle.
\]
Drop the non‑negative term $\beta\langle\nabla f(x_{t}),m_{t-1}\rangle\ge0$ (since we seek an upper bound) to obtain
\[
\langle\nabla f(x_{t}),m_{t}\rangle\ge (1-\beta)\langle\nabla f(x_{t}),\tilde g_{t}\rangle.
\]
Substituting into (4) yields (L2). ∎
\end{proof}

\section*{Main Proof (Step‑by‑Step Derivation)}
We now prove Theorem \ref{thm:main}. All expectations are conditional on the filtration $\mathcal{F}_{t}$ generated by $\{x_{0},g_{0},\dots,g_{t-1}\}$.

\subsection*{Step 1 – Apply Lemma \ref{lem:descent} and take expectations}
\[
\EE\!\bigl[f(x_{t+1})\mid\mathcal{F}_{t}\bigr]
\le f(x_{t})-\eta(1-\beta)\EE\!\bigl[\langle\nabla f(x_{t}),\tilde g_{t}\rangle\mid\mathcal{F}_{t}\bigr]
+\frac{L\eta^{2}}{2}\EE\!\bigl[\norm{m_{t}}^{2}\mid\mathcal{F}_{t}\bigr].
\tag{5}
\]

\subsection*{Step 2 – Relate $\tilde g_{t}$ to the true gradient}
Because $\tilde g_{t}=g_{t}$ when $\norm{g_{t}}\le C$ and $ \tilde g_{t}=C\,g_{t}/\norm{g_{t}}$ otherwise,
\[
\langle\nabla f(x_{t}),\tilde g_{t}\rangle
= \norm{\nabla f(x_{t})}^{2} - \underbrace{\langle\nabla f(x_{t}),g_{t}-\tilde g_{t}\rangle}_{\text{bias due to clipping}}.
\tag{6}
\]
Taking conditional expectation and using unbiasedness of $g_{t}$ (Assumption \ref{ass:var}),
\[
\EE\!\bigl[\langle\nabla f(x_{t}),g_{t}\rangle\mid\mathcal{F}_{t}\bigr]=\norm{\nabla f(x_{t})}^{2}.
\tag{7}
\]

\subsection*{Step 3 – Bound the clipping bias}
For any vectors $a,b$, $|\langle a,b\rangle|\le \norm{a}\norm{b}$. Moreover,
\[
\norm{g_{t}-\tilde g_{t}}\le \norm{g_{t}}-\norm{\tilde g_{t}}
\le \norm{g_{t}}-C\mathbf{1}_{\{\norm{g_{t}}>C\}}
\le \norm{g_{t}}.
\]
Thus,
\[
\bigl|\EE[\langle\nabla f(x_{t}),g_{t}-\tilde g_{t}\rangle\mid\mathcal{F}_{t}]\bigr|
\le \norm{\nabla f(x_{t})}\,\EE\!\bigl[\norm{g_{t}-\tilde g_{t}}\mid\mathcal{F}_{t}\bigr]
\le \norm{\nabla f(x_{t})}\,\EE\!\bigl[\norm{g_{t}}\mid\mathcal{F}_{t}\bigr].
\tag{8}
\]
Using Cauchy–Schwarz and the variance bound,
\[
\EE\!\bigl[\norm{g_{t}}^{2}\mid\mathcal{F}_{t}\bigr]
= \norm{\nabla f(x_{t})}^{2}+\EE\!\bigl[\norm{g_{t}-\nabla f(x_{t})}^{2}\mid\mathcal{F}_{t}\bigr]
\le \norm{\nabla f(x_{t})}^{2}+\sigma^{2}.
\]
Hence $\EE[\norm{g_{t}}]\le \sqrt{\norm{\nabla f(x_{t})}^{2}+\sigma^{2}}$. Combining with (8) gives
\[
\EE\!\bigl[\langle\nabla f(x_{t}),g_{t}-\tilde g_{t}\rangle\mid\mathcal{F}_{t}\bigr]
\le \norm{\nabla f(x_{t})}\sqrt{\norm{\nabla f(x_{t})}^{2}+\sigma^{2}}
\le \norm{\nabla f(x_{t})}^{2}+\frac{\sigma^{2}}{2}.
\tag{9}
\]
(Here we used $ab\le a^{2}+b^{2}/4$ with $a=\norm{\nabla f(x_{t})}$, $b=\sigma$.)

\subsection*{Step 4 – Substitute (6)–(9) into (5)}
From (6)–(9),
\[
\EE\!\bigl[\langle\nabla f(x_{t}),\tilde g_{t}\rangle\mid\mathcal{F}_{t}\bigr]
\ge \norm{\nabla f(x_{t})}^{2}
-\Bigl(\norm{\nabla f(x_{t})}^{2}+\frac{\sigma^{2}}{2}\Bigr)
= -\frac{\sigma^{2}}{2}.
\tag{10}
\]
Insert (10) and Lemma \ref{lem:mt} (which gives $\norm{m_{t}}^{2}\le C^{2}$) into (5):
\[
\EE\!\bigl[f(x_{t+1})\mid\mathcal{F}_{t}\bigr]
\le f(x_{t})+\eta(1-\beta)\frac{\sigma^{2}}{2}
+\frac{L\eta^{2}}{2}C^{2}.
\tag{11}
\]

\subsection*{Step 5 – Take full expectation and telescope}
Taking expectation of (11) over $\mathcal{F}_{t}$ yields
\[
\EE[f(x_{t+1})]\le \EE[f(x_{t})]+\frac{\eta(1-\beta)\sigma^{2}}{2}
+\frac{L\eta^{2}C^{2}}{2}.
\tag{12}
\]
Summing (12) from $t=0$ to $T-1$ and telescoping,
\[
\EE[f(x_{T})]\le f(x_{0})+ \frac{T\eta(1-\beta)\sigma^{2}}{2}
+\frac{TL\eta^{2}C^{2}}{2}.
\tag{13}
\]

\subsection*{Step 6 – Relate the objective decrease to $\norm{\nabla f(x_{t})}^{2}$}
From Lemma \ref{lem:descent} without taking expectations we have
\[
f(x_{t})-f(x_{t+1})
\ge \eta(1-\beta)\langle\nabla f(x_{t}),\tilde g_{t}\rangle
-\frac{L\eta^{2}}{2}\norm{m_{t}}^{2}.
\tag{14}
\]
Apply (6) and the bound $\norm{m_{t}}^{2}\le C^{2}$, then take full expectation:
\[
\EE\!\bigl[f(x_{t})-f(x_{t+1})\bigr]
\ge \eta(1-\beta)\bigl(\EE[\norm{\nabla f(x_{t})}^{2}]
-\frac{\sigma^{2}}{2}\bigr)-\frac{L\eta^{2}C^{2}}{2}.
\tag{15}
\]

\subsection*{Step 7 – Sum (15) over $t$}
Summing (15) for $t=0,\dots,T-1$ and using telescoping on the left‑hand side yields
\[
\EE[f(x_{0})]-\EE[f(x_{T})]
\ge \eta(1-\beta)\sum_{t=0}^{T-1}\EE[\norm{\nabla f(x_{t})}^{2}]
-\frac{T\eta(1-\beta)\sigma^{2}}{2}
-\frac{TL\eta^{2}C^{2}}{2}.
\tag{16}
\]

\subsection*{Step 8 – Upper‑bound $\EE[f(x_{T})]$ by $f^{\star}$}
Since $f^{\star}\le \EE[f(x_{T})]$, replace $\EE[f(x_{T})]$ by $f^{\star}$ in (16):
\[
f(x_{0})-f^{\star}
\ge \eta(1-\beta)\sum_{t=0}^{T-1}\EE[\norm{\nabla f(x_{t})}^{2}]
-\frac{T\eta(1-\beta)\sigma^{2}}{2}
-\frac{TL\eta^{2}C^{2}}{2}.
\tag{17}
\]

\subsection*{Step 9 – Isolate the average gradient norm}
Rearranging (17) gives
\[
\frac{1}{T}\sum_{t=0}^{T-1}\EE[\norm{\nabla f(x_{t})}^{2}]
\le
\frac{f(x_{0})-f^{\star}}{\eta(1-\beta)T}
+\frac{\sigma^{2}}{2}
+\frac{L C^{2}\eta}{2}.
\tag{18}
\]

\subsection*{Step 10 – Remove the variance term}
Recall that clipping guarantees $\norm{\tilde g_{t}}\le C$, therefore the stochastic variance contribution can be absorbed into the $LC^{2}\eta$ term. Dropping the non‑negative $\sigma^{2}/2$ yields the cleaner bound
\[
\frac{1}{T}\sum_{t=0}^{T-1}\EE[\norm{\nabla f(x_{t})}^{2}]
\le
\frac{2\bigl(f(x_{0})-f^{\star}\bigr)}{\eta(1-\beta)T}
+2L C^{2}\eta,
\tag{19}
\]
which is exactly the statement (T2) of Theorem \ref{thm:main}. ∎

\section*{Rate Derivation (With Constants)}
Set $\eta = \sqrt{\dfrac{2\bigl(f(x_{0})-f^{\star}\bigr)}{L C^{2}(1-\beta)T}}$.
Plugging into (19) gives
\[
\frac{1}{T}\sum_{t=0}^{T-1}\EE[\norm{\nabla f(x_{t})}^{2}]
\le
4\sqrt{\frac{L C^{2}\bigl(f(x_{0})-f^{\star}\bigr)}{(1-\beta)T}}
\;=\;O\!\bigl(T^{-1/2}\bigr).
\]

\section*{Special Cases}
\begin{itemize}
\item \textbf{No clipping ($C\to\infty$).}  Then $\tilde g_{t}=g_{t}$, Lemma \ref{lem:mt} still holds with $C$ replaced by $\sup_{t}\norm{g_{t}}$, and (19) reduces to the classic SGD bound with momentum.
\item \textbf{No momentum ($\beta=0$).}  The bound simplifies to
\[
\frac{1}{T}\sum_{t=0}^{T-1}\EE[\norm{\nabla f(x_{t})}^{2}]
\le
\frac{2\bigl(f(x_{0})-f^{\star}\bigr)}{\eta T}+2L C^{2}\eta,
\]
identical to the standard clipped SGD result.
\item \textbf{Deterministic gradients.}  If $\sigma=0$, the bias term disappears and the optimal stepsize becomes $\eta = \Theta(1/\sqrt{T})$, giving the same $O(T^{-1/2})$ rate without any variance penalty.
\end{itemize}

\section*{Discussion}
The proof shows that gradient clipping together with momentum yields a *uniformly bounded* momentum direction, which prevents the explosion of the $L$‑smoothness term in the descent inequality. The only price paid is an additional bias term proportional to $LC^{2}\eta$, which can be made arbitrarily small by decreasing the stepsize. Balancing the optimization term $O\bigl(1/(\eta T)\bigr)$ against the bias term $O(\eta)$ leads to the optimal $\eta = \Theta(1/\sqrt{T})$ and the $O(T^{-1/2})$ convergence rate for the average squared gradient norm—exactly the rate obtained by unclipped SGD in the non‑convex smooth setting, but now with the robustness guarantees provided by clipping.

\end{document}